% Proposed opening for team review, following Simon's scope correction.
% AI-assisted draft. Keep the required disclosure in the assembled report.
\section{Building the model}

When forces act on an elastic object, different parts of it move by
different amounts, changing its shape. How much it deforms depends on its
geometry, the material and where forces and supports are placed. Our aim
is to describe this response for several two-dimensional objects using
the same mathematical framework. The particular shapes and loading cases
will be introduced later.

We describe movement within the plane by a displacement field
$\mathbf{u}(x,y)=(u_x(x,y),u_y(x,y))$. A point initially at $(x,y)$ moves
to $(x+u_x,y+u_y)$, so these two components tell us how far it moves
horizontally and vertically. Moving the whole object by the same amount
does not stretch it. Stretching and changes in angle arise when neighbouring
points move relative to one another; these changes are measured by strain.
The material responds with internal forces, described by stress, which is
force per unit area within the object.

We use small-strain linear elasticity, assuming that displacement
gradients are small and that stress depends linearly on strain. We also
consider static equilibrium: the object is at rest, so the internal forces,
applied loads and forces from the supports must balance. Supports prescribe
some displacement components and prevent the object from moving or
rotating freely as a whole. Their reaction forces are determined as part
of the solution.

The displacement field contains infinitely many unknown values. To obtain
a finite problem, we divide the object's region into triangles and
approximate displacement inside each triangle by linear interpolation of
its values at the three vertices, called nodes. Each node therefore has
two displacement components. Combining the triangles' contributions gives
a stiffness matrix $K$, which describes how nodal displacements produce
internal forces, and a vector $f$ of applied nodal loads. The matrix is
sparse because nodes interact directly only when they share a triangle.

After separating the prescribed displacement components from the unknown
ones, force balance at the free components becomes
\[
  Ax=b.
\]
Here $x$ contains the unknown displacements, $A$ is the corresponding
part of $K$, and $b$ contains the applied loads together with any
contributions from prescribed displacements. With suitable supports and
a stable elastic material, this reduced matrix is symmetric positive
definite. The deformation problem has therefore become a linear algebra
problem. We will study how changing geometry, mesh resolution, loads and
supports affects both the deformation and the numerical solution.
